% ============================================================
% SCHEDE PG — CASA SPINOLA
% @rpg-schema: <https://rpg-schema.org/instances/genova1507/crew-spinola>
%   fitd:hasMember
%     <.../pg-raffaele-spinola>, <.../pg-lucrezia-spinola>,
%     <.../pg-benedetto-adorno>, <.../pg-margherita-gatta>,
%     <.../pg-ippolito-spinola> .
% ============================================================
\chapter{Schede PG — Casa Spinola}
\index{schede PG!Spinola}

% ────── PG 1: RAFFAELE SPINOLA ──────────────────────────────
% @rpg-schema: <.../pg-raffaele-spinola> a fitd:Character ;
%   rdfs:label "Raffaele Spinola il Vecchio" ;
%   fitd:action [ fitd:actionName "Study" ; fitd:dots 4 ] ;
%   fitd:action [ fitd:actionName "Survey" ; fitd:dots 3 ] ;
%   fitd:action [ fitd:actionName "Command" ; fitd:dots 3 ] ;
%   fitd:action [ fitd:actionName "Consort" ; fitd:dots 3 ] ;
%   fitd:action [ fitd:actionName "Sway" ; fitd:dots 4 ] .

\pgheader[SpinolaRed]{Raffaele Spinola il Vecchio}{Il Banchiere / Il Patriarca}{Spinola}
\pgquote{Il denaro non ha odore. Il potere sì — e io ne conosco ogni sfumatura.}

\section*{Background \& Motivazione}
\begin{rulebox}{}
\textbf{Background:} Cinquantotto anni, capo effettivo della famiglia. Raramente visibile in pubblico — opera attraverso intermediari, lettere, e una rete di debitori che copre mezza Europa. Ha finanziato tre guerre, due conclavi, e la costruzione di un palazzo a Roma che porta il nome di qualcun altro. Preferisce così.

\medskip\noindent\textcolor{SpinolaRed}{\textbf{Motivazione:}} La stabilità. Non odia i francesi per principio — li odia perché destabilizzano i mercati e rendono imprevedibili i debiti. Una Genova libera è una Genova dove i contratti vengono rispettati.
\end{rulebox}

\section*{Attributi \& Azioni}

\begin{attributoblock}[SpinolaRed]{INSIGHT — Mente}
  \azione{Caccia}{Hunt}{1}{Trova informazioni tramite la rete}
  \azione{Studio}{Study}{4}{Contabilità, diritto commerciale, codici}
  \azione{Osserva}{Survey}{3}{Legge le persone come legge un bilancio}
  \azione{Ingegno}{Tinker}{2}{Sistemi, strutture organizzative}
\end{attributoblock}
\begin{attributoblock}[SpinolaRed]{PROWESS — Corpo}
  \azione{Finezza}{Finesse}{1}{Scrittura, documenti, sigilli}
  \azione{Ombra}{Prowl}{0}{Non se ne ha bisogno quando si ha il potere}
  \azione{Combattimento}{Skirmish}{0}{Non da decenni}
  \azione{Forza bruta}{Wreck}{0}{Assolutamente no}
\end{attributoblock}

\begin{attributoblock}[SpinolaRed]{RESOLVE — Spirito}

  \azione{Percezione}{Attune}{2}{Sente l'equilibrio di una stanza}
  \azione{Comando}{Command}{3}{Non alza mai la voce. Non ne ha bisogno.}
  \azione{Relazioni}{Consort}{3}{Conosce personalmente chiunque conti qualcosa}
  \azione{Persuasione}{Sway}{4}{Negozia, convince, compra e vende consenso}

\end{attributoblock}

\medskip
\begin{pgclockbox}{Stress \hfill \clock{9}{0}}
\small\textit{Traumi: $\square$ Freddo \quad $\square$ Duro \quad $\square$ Ossessionato \quad $\square$ Paranoico \quad $\square$ Irrazionale \quad $\square$ Spezzato}
\end{pgclockbox}

\section*{Abilità Speciale}
\begin{spinolabox}{Il Registro dei Debiti}
Raffaele può dichiarare che qualsiasi PNG ha un debito con la famiglia Spinola. Il PNG è obbligato a rispettare una richiesta ragionevole. Una richiesta irragionevole richiede un Tiro su Persuasione da Posizione Controllata. \textit{Tre volte per sessione. Solo sui PNG.}
\end{spinolabox}

\section*{Legami}
\begin{table}[h]\centering\renewcommand{\arraystretch}{1.4}
\begin{tabularx}{\linewidth}{@{}L{2.6cm} X@{}}
\toprule\rowcolor{SpinolaRed}
\textcolor{white}{\textbf{PG}} & \textcolor{white}{\textbf{Legame}} \\
\midrule
Lucrezia & La nostra intermediaria. Brava. Forse troppo. Ultimamente le domande vanno oltre il mandato. \\
\rowcolor{SpinolaRed!8}
Ser Benedetto & Il notaio sa tutto. L'ho scelto perché non ha alternative alla lealtà. Ma la lealtà forzata è fragile. \\
Margherita & La mia agente di campo. Non la controllo — la oriento. C'è differenza. \\
\rowcolor{SpinolaRed!8}
Ippolito & Mio nipote. Brillante e impaziente. Lo amo e ho paura di lui. \\
\bottomrule
\end{tabularx}
\end{table}

\section*{Equipaggiamento}
\begin{goldlist}
\item Il Registro: libro mastro cifrato con i debiti di trecento persone influenti
\item Sigillo in ceralacca degli Spinola — un documento firmato con questo vale oro
\item Lettere di credito su tre banche europee — liquidità immediata ovunque
\item Antidoto generico per i veleni più comuni. Lo porta sempre. Da trent'anni.
\end{goldlist}

\clearpage
\begin{secretbox}
\textbf{Segreto:} Raffaele ha già deciso che se la rivolta fallisce, negozierà con i francesi. Ha preparato una memoria per Luigi XII in cui offre i servizi del Banco in cambio dell'immunità. Il documento esiste. Nessun altro Spinola lo sa.

\medskip\noindent\textcolor{SecretRed}{\textbf{Agenda:}} \textit{Controllare l'esito della rivolta. Il banco Spinola deve sopravvivere — qualunque sia il vincitore.}
\end{secretbox}

\clearpage

% ────── PG 2: LUCREZIA ──────────────────────────────────────
% @rpg-schema: <.../pg-lucrezia-spinola> a fitd:Character ;
%   rdfs:label "Lucrezia d'Oria in Spinola" ;
%   fitd:action [ fitd:actionName "Study" ; fitd:dots 3 ] ;
%   fitd:action [ fitd:actionName "Consort" ; fitd:dots 4 ] ;
%   fitd:action [ fitd:actionName "Sway" ; fitd:dots 3 ] .

\pgheader[SpinolaRed]{Lucrezia d'Oria in Spinola}{La Diplomatica / L'Agente Estero}{Spinola}
\pgquote{I segreti che porto non mi appartengono. È per questo che posso tenerli.}

\section*{Background \& Motivazione}
\begin{rulebox}{}
\textbf{Background:} Trentadue anni, nata Doria — un matrimonio di convenienza trasformato in arma. Intermediaria tra gli Spinola e le corti straniere, si muove tra Genova, Ferrara e Milano. È ufficialmente «in visita per la salute». Tutti sanno che è una bugia. Nessuno sa la verità.

\medskip\noindent\textcolor{SpinolaRed}{\textbf{Motivazione:}} La lealtà la spaventa: ne ha troppa verso troppe persone. Vuole che la rivolta risolva il conflitto così da poter scegliere finalmente da che parte stare.
\end{rulebox}

\section*{Attributi \& Azioni}

\begin{attributoblock}[SpinolaRed]{INSIGHT — Mente}
  \azione{Caccia}{Hunt}{2}{Trova info nei canali diplomatici}
  \azione{Studio}{Study}{3}{Lingue, protocollo di corte, codici cifrati}
  \azione{Osserva}{Survey}{2}{Legge le corti, le dinamiche sociali}
  \azione{Ingegno}{Tinker}{1}{Codici, steganografia}
\end{attributoblock}
\begin{attributoblock}[SpinolaRed]{PROWESS — Corpo}
  \azione{Finezza}{Finesse}{2}{Movimenti sociali precisi, gestualità calibrata}
  \azione{Ombra}{Prowl}{1}{Si muove negli ambienti aristocratici invisibilmente}
  \azione{Combattimento}{Skirmish}{0}{Nessuna esperienza}
  \azione{Forza bruta}{Wreck}{0}{No}
\end{attributoblock}

\begin{attributoblock}[SpinolaRed]{RESOLVE — Spirito}

  \azione{Percezione}{Attune}{2}{Sente le correnti emotive di una stanza}
  \azione{Comando}{Command}{1}{Autorità discreta, non imposta}
  \azione{Relazioni}{Consort}{4}{A casa in qualsiasi corte europea}
  \azione{Persuasione}{Sway}{3}{Convince attraverso il fascino, non la logica}

\end{attributoblock}

\medskip
\begin{pgclockbox}{Stress \hfill \clock{9}{0}}
\small\textit{Traumi: $\square$ Freddo \quad $\square$ Duro \quad $\square$ Ossessionato \quad $\square$ Paranoico \quad $\square$ Irrazionale \quad $\square$ Spezzato}
\end{pgclockbox}

\section*{Abilità Speciale}
\begin{spinolabox}{Passaporto Diplomatico}
Lucrezia può muoversi in qualsiasi ambiente sociale — incluse aree controllate dai francesi o palazzi delle altre famiglie — senza essere fermata, purché non compia azioni palesemente ostili. Include portare con sé un'altra persona. \textit{Sempre attiva. Decade nella scena se compie azioni ostili.}
\end{spinolabox}

\section*{Legami}
\begin{table}[h]\centering\renewcommand{\arraystretch}{1.4}
\begin{tabularx}{\linewidth}{@{}L{2.6cm} X@{}}
\toprule\rowcolor{SpinolaRed}
\textcolor{white}{\textbf{PG}} & \textcolor{white}{\textbf{Legame}} \\
\midrule
Raffaele & Lo rispetto. La memoria che ha preparato per i francesi — l'ho trovata. Non sa che l'ho letta. \\
\rowcolor{SpinolaRed!8}
Ser Benedetto & Sa troppo. Lo compatisco — essere il custode dei segreti altrui è una prigione. \\
Margherita & Due donne che sopravvivono in mondi fatti per escluderle. Diffidenza reciproca. \\
\rowcolor{SpinolaRed!8}
Ippolito & È innamorato di me. Lo sa. Lo lascio credere di non saperlo. Mi vergogno di questo. \\
\bottomrule
\end{tabularx}
\end{table}

\section*{Equipaggiamento}
\begin{goldlist}
\item Salvacondotto firmato dal governatore francese (ottenuto prima della crisi)
\item Tre lettere sigillate indirizzate a corti diverse — da consegnare solo in emergenza
\item Gioiello con vano nascosto contenente un documento in codice
\item Lista dei debiti personali del governatore con la famiglia Spinola — la sua copia
\end{goldlist}

\clearpage
\begin{secretbox}
\textbf{Segreto:} Lucrezia ha trovato la memoria segreta di Raffaele per i francesi. Non sa ancora cosa farne. Intanto la tiene. Questo la rende pericolosa per Raffaele — anche se lui non lo sa ancora.

\medskip\noindent\textcolor{SecretRed}{\textbf{Agenda:}} \textit{Capire se la lealtà agli Spinola vale il costo che sta pagando. Se trova evidenza che Raffaele è disposto a sacrificarla, uscirà dalla famiglia con tutto quello che sa.}
\end{secretbox}

\clearpage

% ────── PG 3: BENEDETTO ADORNO ──────────────────────────────
% @rpg-schema: <.../pg-benedetto-adorno> a fitd:Character ;
%   rdfs:label "Ser Benedetto Adorno" ;
%   fitd:action [ fitd:actionName "Study" ; fitd:dots 4 ] ;
%   fitd:action [ fitd:actionName "Finesse" ; fitd:dots 3 ] ;
%   fitd:action [ fitd:actionName "Tinker" ; fitd:dots 3 ] .

\pgheader[SpinolaRed]{Ser Benedetto Adorno}{Il Notaio / Il Custode dei Segreti}{Spinola}
\pgquote{Ho scritto contratti che hanno cambiato il destino di città. Nessuno sa il mio nome.}

\section*{Background \& Motivazione}
\begin{rulebox}{}
\textbf{Background:} Cinquantadue anni, notaio del Palazzo di San Giorgio, uomo di fiducia degli Spinola da vent'anni. Custodisce i veri libri contabili del banco. È un uomo profondamente onesto in un ruolo che richiede assoluta disonestà verso il mondo esterno. Questa contraddizione lo sta consumando.

\medskip\noindent\textcolor{SpinolaRed}{\textbf{Motivazione:}} Vuole che la rivolta finisca in modo da poter smettere di mentire. Una Genova con leggi nuove potrebbe permettergli di rendere pubblici i conti veri senza rischiare la vita.
\end{rulebox}

\section*{Attributi \& Azioni}

\begin{attributoblock}[SpinolaRed]{INSIGHT — Mente}
  \azione{Caccia}{Hunt}{0}{Non è nel suo repertorio}
  \azione{Studio}{Study}{4}{Diritto, contabilità, paleografia, archivistica}
  \azione{Osserva}{Survey}{2}{Nota dettagli nei documenti che altri ignorano}
  \azione{Ingegno}{Tinker}{3}{Codici, sigilli, falsificazione (sa farlo)}
\end{attributoblock}
\begin{attributoblock}[SpinolaRed]{PROWESS — Corpo}
  \azione{Finezza}{Finesse}{3}{Scrittura, documenti, copie perfette}
  \azione{Ombra}{Prowl}{0}{Non è la sua specialità}
  \azione{Combattimento}{Skirmish}{0}{Mai}
  \azione{Forza bruta}{Wreck}{0}{No}
\end{attributoblock}

\begin{attributoblock}[SpinolaRed]{RESOLVE — Spirito}

  \azione{Percezione}{Attune}{1}{Sente quando una firma è falsa}
  \azione{Comando}{Command}{0}{Non ha autorità. È un funzionario.}
  \azione{Relazioni}{Consort}{2}{Notai, funzionari, burocrazia}
  \azione{Persuasione}{Sway}{2}{Convince con la logica e i documenti}

\end{attributoblock}

\medskip
\begin{pgclockbox}{Stress \hfill \clock{9}{0}}
\small\textit{Traumi: $\square$ Freddo \quad $\square$ Duro \quad $\square$ Ossessionato \quad $\square$ Paranoico \quad $\square$ Irrazionale \quad $\square$ Spezzato}
\end{pgclockbox}

\section*{Abilità Speciale}
\begin{spinolabox}{Il Documento Giusto}
Ser Benedetto può produrre o modificare qualsiasi documento in modo convincente — contratto, salvacondotto, lettera di credito, mandato di arresto. Il documento è sempre legalmente plausibile. Documenti semplici: nessun tiro. Documenti complessi con sigilli ufficiali: Tiro su Studio.
\end{spinolabox}

\section*{Legami}
\begin{table}[h]\centering\renewcommand{\arraystretch}{1.4}
\begin{tabularx}{\linewidth}{@{}L{2.6cm} X@{}}
\toprule\rowcolor{SpinolaRed}
\textcolor{white}{\textbf{PG}} & \textcolor{white}{\textbf{Legame}} \\
\midrule
Raffaele & Mi ha salvato vent'anni fa. Gli devo la vita — e questa è la catena più solida che esista. \\
\rowcolor{SpinolaRed!8}
Lucrezia & È troppo intelligente. Mi guarda come se sapesse cose. Forse ne sa. \\
Margherita & Non capisco come funzioni. Eppure ci troviamo sempre a condividere le stesse informazioni. \\
\rowcolor{SpinolaRed!8}
Ippolito & Mi spaventa. Ha l'impazienza di chi non capisce che certe cose si rompono se tirate troppo. \\
\bottomrule
\end{tabularx}
\end{table}

\section*{Equipaggiamento}
\begin{goldlist}
\item Copia dei libri contabili segreti del banco (criptata)
\item Sigillo ufficiale del Palazzo di San Giorgio — può autenticare documenti istituzionali
\item Lista dei conti falsi usati per il Colpo Riservato — da distruggere dopo
\item Lettera anonima pronta per l'arcivescovo con info sul banco. Non l'ha ancora inviata.
\end{goldlist}

\clearpage
\begin{secretbox}
\textbf{Segreto:} Benedetto ha un figlio illegittimo di diciott'anni che gli Spinola mantengono in segreto. Il ragazzo ha recentemente chiesto di incontrarsi. Se questo avviene durante la sessione, Benedetto sarà distratto nel momento peggiore.

\medskip\noindent\textcolor{SecretRed}{\textbf{Agenda:}} \textit{Proteggere il figlio. Se la rivolta crea le condizioni per riconoscerlo pubblicamente senza rischi, Benedetto firmerà qualsiasi cosa, per chiunque.}
\end{secretbox}

\clearpage

% ────── PG 4: MARGHERITA LA GATTA ───────────────────────────
% @rpg-schema: <.../pg-margherita-gatta> a fitd:Character ;
%   rdfs:label "Margherita 'la Gatta'" ;
%   fitd:action [ fitd:actionName "Hunt" ; fitd:dots 3 ] ;
%   fitd:action [ fitd:actionName "Survey" ; fitd:dots 4 ] ;
%   fitd:action [ fitd:actionName "Prowl" ; fitd:dots 4 ] ;
%   fitd:action [ fitd:actionName "Consort" ; fitd:dots 3 ] .

\pgheader[SpinolaRed]{Margherita «la Gatta»}{L'Agente di Campo / L'Informatrice}{Spinola}
\pgquote{Non sono una spia. Sono una donna che presta attenzione. La differenza la fanno gli altri.}

\section*{Background \& Motivazione}
\begin{rulebox}{}
\textbf{Background:} Trentanove anni, origini nei bassifondi di Soziglia. È entrata nella rete Spinola per caso — ha sentito qualcosa che non avrebbe dovuto sentire, e invece di essere eliminata è stata reclutata. È la loro agente più capace sul campo.

\medskip\noindent\textcolor{SpinolaRed}{\textbf{Motivazione:}} Un appartamento suo, abbastanza denaro da scegliere, e nessuno che la controlli. La rivolta è l'occasione per passare dalla sopravvivenza alla libertà.
\end{rulebox}

\section*{Attributi \& Azioni}

\begin{attributoblock}[SpinolaRed]{INSIGHT — Mente}
  \azione{Caccia}{Hunt}{3}{Pedina, localizza, segue — chiunque, ovunque}
  \azione{Studio}{Study}{1}{Sa leggere e scrivere — non ovvio per la sua origine}
  \azione{Osserva}{Survey}{4}{La sua sopravvivenza ha sempre dipeso da questo}
  \azione{Ingegno}{Tinker}{1}{Sa aprire serrature semplici}
\end{attributoblock}
\begin{attributoblock}[SpinolaRed]{PROWESS — Corpo}
  \azione{Finezza}{Finesse}{2}{Borseggio, movimenti precisi in spazi affollati}
  \azione{Ombra}{Prowl}{4}{Come il suo soprannome. Sparisce e riappare.}
  \azione{Combattimento}{Skirmish}{2}{Coltello, qualsiasi oggetto improvvisato}
  \azione{Forza bruta}{Wreck}{0}{Non è la strategia giusta per lei}
\end{attributoblock}

\begin{attributoblock}[SpinolaRed]{RESOLVE — Spirito}

  \azione{Percezione}{Attune}{2}{Sente quando qualcosa non va. Raramente sbaglia.}
  \azione{Comando}{Command}{0}{Non comanda — suggerisce}
  \azione{Relazioni}{Consort}{3}{Taverne, mercati, reti femminili dei quartieri}
  \azione{Persuasione}{Sway}{2}{Convince con la verità più che con le parole}

\end{attributoblock}

\medskip
\begin{pgclockbox}{Stress \hfill \clock{9}{0}}
\small\textit{Traumi: $\square$ Freddo \quad $\square$ Duro \quad $\square$ Ossessionato \quad $\square$ Paranoico \quad $\square$ Irrazionale \quad $\square$ Spezzato}
\end{pgclockbox}

\section*{Abilità Speciale}
\begin{spinolabox}{Invisibile ai Potenti}
Margherita non viene percepita come minaccia da nessuno con status sociale superiore al suo (mercanti, nobili, ufficiali francesi). Può muoversi liberamente in ambienti in cui sarebbe in teoria fuori posto. \textit{Sempre attiva. Decade nella scena se attira attenzione su di sé.}
\end{spinolabox}

\section*{Legami}
\begin{table}[h]\centering\renewcommand{\arraystretch}{1.4}
\begin{tabularx}{\linewidth}{@{}L{2.6cm} X@{}}
\toprule\rowcolor{SpinolaRed}
\textcolor{white}{\textbf{PG}} & \textcolor{white}{\textbf{Legame}} \\
\midrule
Raffaele & Mi rispetta perché sono utile. Non sono sicura di come la prenderei se smettessi di esserlo. \\
\rowcolor{SpinolaRed!8}
Lucrezia & Due donne che sopravvivono in mondi fatti per escluderle. Potremmo essere alleate o nemiche. \\
Ser Benedetto & Il notaio è sul punto di rompersi. Lo osservo. Non so ancora se intervenire o aspettare. \\
\rowcolor{SpinolaRed!8}
Ippolito & L'ho salvato da un'imboscata. Da allora mi segue come un cagnolino. Scomodo e strano. \\
\bottomrule
\end{tabularx}
\end{table}

\section*{Equipaggiamento}
\begin{goldlist}
\item Cambio di abiti completo nascosto in tre posti diversi nei caruggi
\item Piccolo coltello con lama ceramica — quasi impossibile da rilevare
\item Rete di quattro informatori nei sestieri principali (chiamati «le mie sorelle»)
\item Chiave che apre la porta secondaria del Palazzo di San Giorgio — non sa come ce l'ha
\end{goldlist}

\clearpage
\begin{secretbox}
\textbf{Segreto:} Margherita ha scoperto che uno degli informatori di Raffaele è in realtà un agente dei Fieschi. Lo sa da due settimane e non lo ha riferito. Sta cercando di capire se può usare questo canale per i propri scopi.

\medskip\noindent\textcolor{SecretRed}{\textbf{Agenda:}} \textit{Accumulare abbastanza informazioni da poter negoziare la propria libertà con chiunque vinca. Alla fine della sessione, Margherita vuole poter scomparire — con abbastanza oro e abbastanza segreti da iniziare altrove.}
\end{secretbox}

\clearpage

% ────── PG 5: IPPOLITO SPINOLA ───────────────────────────────
% @rpg-schema: <.../pg-ippolito-spinola> a fitd:Character ;
%   rdfs:label "Ippolito Spinola" ;
%   fitd:action [ fitd:actionName "Study" ; fitd:dots 3 ] ;
%   fitd:action [ fitd:actionName "Command" ; fitd:dots 3 ] ;
%   fitd:action [ fitd:actionName "Sway" ; fitd:dots 3 ] .

\pgheader[SpinolaRed]{Ippolito Spinola}{Il Giovane Erede / Il Giocatore d'Azzardo}{Spinola}
\pgquote{Mio zio pensa che io sia impaziente. Ha ragione. Ma la pazienza è solo la paura travestita.}

\section*{Background \& Motivazione}
\begin{rulebox}{}
\textbf{Background:} Ventidue anni, nipote di Raffaele, destinato a ereditare la guida del banco. È brillante — forse il più brillante della famiglia — e lo sa nel modo peggiore possibile. Cinque anni a Bologna, tre a osservare Raffaele, zero fallimenti significativi. Questo è il suo punto cieco.

\medskip\noindent\textcolor{SpinolaRed}{\textbf{Motivazione:}} Dimostrare a Raffaele che è pronto. Non aspettare. Non fare il secondo. Essere lui il protagonista della liberazione di Genova.
\end{rulebox}

\section*{Attributi \& Azioni}

\begin{attributoblock}[SpinolaRed]{INSIGHT — Mente}
  \azione{Caccia}{Hunt}{2}{Trova vulnerabilità nei sistemi e nelle persone}
  \azione{Studio}{Study}{3}{Diritto, filosofia, matematica — formazione universitaria}
  \azione{Osserva}{Survey}{2}{Intelligente ma vede a volte ciò che vuole vedere}
  \azione{Ingegno}{Tinker}{2}{Schemi, piani, strutture — ama costruire sistemi}
\end{attributoblock}
\begin{attributoblock}[SpinolaRed]{PROWESS — Corpo}
  \azione{Finezza}{Finesse}{2}{Scherma, eleganza fisica — si è allenato}
  \azione{Ombra}{Prowl}{1}{Discreto quando vuole}
  \azione{Combattimento}{Skirmish}{2}{Duellist — addestrato, non ancora provato}
  \azione{Forza bruta}{Wreck}{1}{Ci prova. A volte funziona.}
\end{attributoblock}

\begin{attributoblock}[SpinolaRed]{RESOLVE — Spirito}

  \azione{Percezione}{Attune}{1}{Intuisce le emozioni attraverso la logica}
  \azione{Comando}{Command}{3}{Vuole comandare. Riesce più spesso di quanto dovrebbe.}
  \azione{Relazioni}{Consort}{2}{Circoli universitari, giovani nobili, taverne eleganti}
  \azione{Persuasione}{Sway}{3}{Entusiasmo genuino come strumento persuasivo}

\end{attributoblock}

\medskip
\begin{pgclockbox}{Stress \hfill \clock{9}{0}}
\small\textit{Traumi: $\square$ Freddo \quad $\square$ Duro \quad $\square$ Ossessionato \quad $\square$ Paranoico \quad $\square$ Irrazionale \quad $\square$ Spezzato}
\end{pgclockbox}

\section*{Abilità Speciale}
\begin{spinolabox}{Il Colpo Audace}
Quando Ippolito propone un piano rischioso e almeno un altro PG del tavolo lo segue, entrambi tirano con +1d. Se il piano fallisce con Conseguenze Gravi, Ippolito assorbe metà delle conseguenze al posto del compagno. \textit{Una volta per scena. Il piano deve essere genuinamente rischioso — il GM decide.}
\end{spinolabox}

\section*{Legami}
\begin{table}[h]\centering\renewcommand{\arraystretch}{1.4}
\begin{tabularx}{\linewidth}{@{}L{2.6cm} X@{}}
\toprule\rowcolor{SpinolaRed}
\textcolor{white}{\textbf{PG}} & \textcolor{white}{\textbf{Legame}} \\
\midrule
Raffaele & Lo amo come un padre. Lo temo come un dio. Voglio che sia orgoglioso. Questa è la mia debolezza. \\
\rowcolor{SpinolaRed!8}
Lucrezia & Sono innamorato in modo stupido e ovvio. Lei finge di non saperlo. Entrambi fingiamo. \\
Ser Benedetto & L'unico che mi tratta come un adulto senza flattarmi. Gli voglio bene per questo. \\
\rowcolor{SpinolaRed!8}
Margherita & Mi ha salvato la vita. Questo crea un tipo di legame che non so ancora come gestire. \\
\bottomrule
\end{tabularx}
\end{table}

\section*{Equipaggiamento}
\begin{goldlist}
\item Spada da lato di ottima fattura — regalo di Raffaele per la laurea
\item Taccuino con il piano alternativo per il Colpo Riservato — non autorizzato
\item Lettera di Lucrezia che non avrebbe dovuto leggere ma ha letto
\item Denaro contante — più di quanto un giovane della sua età dovrebbe avere in tasca
\end{goldlist}

\clearpage
\begin{secretbox}
\textbf{Segreto:} Ippolito ha già contattato il rappresentante francese proponendosi come intermediario per una trattativa separata. Non per tradire la famiglia — per dimostrare di poter chiudere affari da solo. Il suo appuntamento di risposta è stasera.

\medskip\noindent\textcolor{SecretRed}{\textbf{Agenda:}} \textit{Compiere almeno una cosa di cui la famiglia possa essere orgogliosa — senza il permesso di Raffaele. Per dimostrare di esistere come agente autonomo, non come promessa di futuro.}
\end{secretbox}
